%This is a template file for use of iopjournal.cls

\documentclass{iopjournal}

% Options
% 	[anonymous]	Provides output without author names, affiliations or acknowledgments to facilitate double-anonymous peer-review

\begin{document}

\articletype{Article type} %	 e.g. Paper, Letter, Topical Review...

\title{Interruption Success or Failure Caused by Arc Behavior in Air Switchgear}

\author{Shuto Igarashi$^1$ , Akio Saitsu$^2$ , Naoya Yamaguchi$^2$ , Yuki Inada$^{1,*}$}

\affil{$^1$Division of Mathematics, Electronics and Information Sciences, Saitama University, Shimo-Okubo, Sakura-ku, Saitama, Saitama 338-8570, Japan}

\affil{$^2$Togami Electric Mfg. Co., Ltd, Saga, Japan}

\affil{$^*$Author to whom any correspondence should be addressed.}

\email{s.igarashi.232@ms.saitama-u.ac.jp and inada@mail.saitama-u.ac.jp}

\keywords{air switchgear, arc behavior , spectroscopy}


\begin{abstract}
To clarify the factors leading to arc reignition in air switchgear, the spatiotemporal variations of arc behavior, temperature, 
and electrical conductivity, as well as hot gas emissions near the first current zero, were measured. 
In cases with short interruption times, the arc spot on the movable electrode tended to move away from the contact gap, 
and the electrical conductivity near the spot ranged from $10^{-1}$ to $10^{2}$~S/m, approximately one order of magnitude lower than in unsuccessful cases. 
This suggests that the spot position immediately before the current zero is a critical factor for successful thermal interruption. 
In contrast, in cases where reignition occurred, the emission intensity of gas ejected from the arc-quenching chamber was weak, 
and the ratio of the total gas to the total arc path temperature was less than half that in successful interruptions. 
Since the injected electrical energy between the electrodes showed no significant difference between successful and unsuccessful cases, 
it is suggested that the strong gas emission contributes to suppressing arc path luminosity, thereby preventing reignition.
\end{abstract}

\section{Section title}
Sample text inserted for demonstration. Organize the main text of your article using section headings, and include any equations, figures, tables, lists etc using your preferred \LaTeX\ packages and commands. Example code for a figure and a table is given below, but you do not have to use this format. For general guidance on using \LaTeX , including information on figures, tables, equations and references, please refer to documents such as the \LaTeX\ WikiBook: \href{https://en.wikibooks.org/wiki/LaTeX}{https://en.wikibooks.org/wiki/LaTeX}.

Note that clarity of presentation is the most important consideration when preparing your article for submission. It is not necessary to format your article in the style used for published articles in the journal.

\subsection{Subsection title}
Sample text inserted for demonstration, including links to figure \ref{fig1} and table \ref{tab1}.

\subsubsection{Subsubsection heading}
Sample text inserted for demonstration.

% \begin{figure}
%  \centering
%         \includegraphics[width=0.5\textwidth]{figure1}
%  \caption{Text describing the figure and the main conclusions drawn from it. To make your figures accessible to as many readers as possible, try to avoid using colour as the only means of conveying information. For example, in charts and graphs use different line styles and symbols. Further information is available in the online guide: \href{https://publishingsupport.iopscience.iop.org/publishing-support/authors/authoring-for-journals/writing-journal-article/\#figures}{https://publishingsupport.iopscience.iop.org/publishing-support/authors/authoring-for-journals/writing-journal-article/\#figures}}
% \label{fig1}
% \end{figure}


\begin{table}
\caption{Caption text describing the table. Adapt the template table below or replace with a new table. To add more tables, copy and paste the whole {\tt \textbackslash begin\{table\}...\textbackslash end\{table\}} block.}
\centering
\begin{tabular}{l c c c}
\hline
Column heading & Column heading & Column heading & Column heading \\
\hline
Data row 1 & 1.0 & 1.5 & 2.0 \\
Data row 2 & 2.0 & 2.5 & 3.0 \\
Data row 3 & 3.0 & 3.5 & 4.0 \\
\hline
\end{tabular}
\label{tab1}
\end{table}

%
% Each of the commands below will create an unnumbered section with the appropriate heading.
% Remove any sections that are not relevant for your article.
% All sections except suppdata will be removed if the [anonymous] option is used.
% See iopjournal-guidelines.pdf for more information.
%

\ack{Sample text inserted for demonstration.}

\funding{Sample text inserted for demonstration.}
% This section is a list of funder names and grant numbers

\roles{Sample text inserted for demonstration.}
% List author names and the contributions made to the article, using terms from the NISO Contributor Roles Taxonomy (CRediT) https://credit.niso.org

\data{Sample text inserted for demonstration.}
% For more information on IOP Publishing's research data policy see: https://publishingsupport.iopscience.iop.org/questions/research-data/

\suppdata{Sample text inserted for demonstration.}


\section*{References}



\end{document}


